\chapter{Business case}\label{ch:business}

\section{Value proposition}

\begin{quote}
\emph{Other teams map what you say. We find you when you don't know where you are, by asking the right question.}
\end{quote}

For a military medical chain, Near does three things that save time when time decides who lives:
\begin{itemize}
  \item it gives a position, with an honest radius, for casualties whose reporters can't give one;
  \item it turns a free-form spoken report into a MIST card, timers and an evacuation order without a second person typing;
  \item it plans who goes where, and by which way, around the threats that other units just reported.
\end{itemize}
It does this with what units already carry: a radio, a compass, sometimes a phone, and one laptop at the
collection point.

\nanofig{segments}{Who Near is for: military medical chains first, then civilian rescue where people can't say where they are.}
{Clean business infographic, 16:9, plain white background, flat vector, navy text, cyan and amber accents. A central circle labelled
``Near: locate from what people see''. Around it, five segment cards connected by thin lines: ``Military medical services
(CASEVAC, CCP)'' with a stretcher icon; ``Command posts and training centres'' with a tent and radio icon; ``Ukraine and frontline
partners'' with a radio and drone-crossed icon; ``Emergency call centres (lost callers)'' with a headset icon; ``Mountain rescue and
disaster response'' with a mountain and helicopter icon. Primary segments (first three) in cyan, secondary (last two) in amber, small
legend ``primary'' / ``next''. No flags, no logos. Exactly these labels. Colour palette (use only these): navy \#0D1110, cyan \#1AA7C4, amber \#E09A2B, red \#E5484D, green \#3FA877, grey \#8B9A94, plus white for text on dark shapes.}

\section{Market context}

Analyst estimates vary widely because each firm draws the market's boundaries differently. We quote them only to
show the order of magnitude of the budgets Near plugs into. Our own sizing is bottom-up (next section).

\begin{table}[htbp]\centering\small
\caption{Published market estimates (vendor research, 2025 values).}\label{tab:market}
\begin{tabularx}{\linewidth}{Xrl}
\toprule
Market & 2025 size & Source\\
\midrule
Tactical communication (broad) & USD 26.2 bn, $\sim$7 \%/yr & Fortune Business Insights \cite{fbi-tac}\\
Military tactical radios (narrow) & USD 5.88 bn, 4.6 \%/yr & Fortune Business Insights \cite{fbi-radio}\\
Military medicine (combat casualty care to aeromedical evacuation) & USD 12.5 bn, 6.8 \%/yr & TrendX Insights \cite{milmed}\\
Battlefield medical equipment & USD 8.5 bn, 8.1 \%/yr & SNS Insider \cite{battlemed}\\
\bottomrule
\end{tabularx}
\end{table}

Public money is moving the same way. The European Defence Fund has about €7.3--7.95 bn for 2021--2027, and Ukraine is an
associated country that can join EDF consortia \cite{edf,consilium}. National innovation agencies run defence
programmes for start-ups, among them the French AID, NATO DIANA and Ukraine's Brave1 cluster.

Demand is driven by the two facts in Chapter~\ref{ch:problem}. GPS interference is now routine over Europe (about 123,000
flights affected in four months of 2025 around the Baltic \cite{gps123k}), and drones have stretched evacuation times
from an hour to a day or more \cite{goldenday,wotr}.

\section{Bottom-up sizing method}

We size the opportunity by the number of places where Near would run, not by a share of a vendor estimate:
\begin{equation}
  \text{Revenue} = N_{\text{sites}} \times P_{\text{site}} + N_{\text{areas}} \times P_{\text{area}} + \text{training and support}.
\end{equation}
\begin{itemize}
  \item $N_{\text{sites}}$: casualty collection points, battalion aid stations and command posts that receive casualty reports;
        on the civil side, dispatch centres and rescue bases.
  \item $P_{\text{site}}$: a yearly licence per laptop running Near. Phones and radios add no cost.
  \item $N_{\text{areas}}$, $P_{\text{area}}$: map data packs (landmarks, buildings, LiDAR, paths, tiles, precomputed line of
        sight) per operating area, rebuilt when the area changes.
\end{itemize}
These parameters are hypotheses to test with the first customers; Section~\ref{sec:validate} lists how. We deliberately
give no revenue forecast in this report: a number without a pilot behind it would be made up.

\section{Business model}

\begin{table}[htbp]\centering\small
\caption{Revenue lines (hypotheses to validate).}\label{tab:revenue}
\begin{tabularx}{\linewidth}{>{\raggedright}p{3.6cm}X>{\raggedright\arraybackslash}p{3.6cm}}
\toprule
Line & What the customer gets & Buyer\\
\midrule
Site licence (yearly) & Near on one CCP or command-post laptop, updates, offline models & Military medical service, unit, rescue base\\
Area data pack & An operating area prepared from open and national data: landmarks, buildings, LiDAR line of sight, walking network, offline tiles & Same, per area and refresh\\
Training & Half-day course: spoken bearings, radio procedure with Near, the offline phone & Training centres\\
Integration & Packets and positions into existing systems (battle management, ATAK-type tools), custom radios & Integrators, primes\\
Funded R\&D & Pilots and adaptations paid by defence innovation programmes & EDF consortia, AID, DIANA, Brave1\\
\bottomrule
\end{tabularx}
\end{table}

\section{Why it scales: mass manufacturability}

\begin{itemize}
  \item \textbf{No hardware to build.} Near is software on a commercial laptop. Our demo machine is a gaming laptop with an
        8 GB GPU, and a CPU-only laptop runs it with the smaller Whisper. The field side uses radios and compasses units already
        own; the phone is optional.
  \item \textbf{Near-zero marginal cost.} A new site is a copy of the software and its models (about 2 GB).
  \item \textbf{Areas are built from data, not by hand.} Every area input is open or national data: OpenStreetMap, IGN BD TOPO
        and LiDAR HD in France, Copernicus elevation elsewhere. The scripts fetch and precompute an area in minutes.
  \item \textbf{Robust supply chain.} No custom chips, no proprietary cloud, permissive model licences (MIT, Apache 2.0,
        BSD-3) and our own trained model.
\end{itemize}

\section{Competition}

\begin{table}[htbp]\centering\small
\caption{Competitive landscape.}\label{tab:competition}
\begin{tabularx}{\linewidth}{>{\raggedright}p{3.4cm}XX}
\toprule
Alternative & Strength & Gap Near fills\\
\midrule
ATAK-type situational awareness apps and battle management systems & Standard, widely fielded, rich maps & Positions come from device GPS; no help when the reporter is unsure or jammed\\
MEDEVAC request apps, digital casualty cards & Structured medical data & Need a known location; typing under fire\\
Voice-to-map solutions (other challenge teams) & Fast speech-to-pin, geocoding, trust scores & Assume the speaker can name his position\\
Acoustic gunshot / drone sensors & Accurate source location & Dedicated hardware, fixed sites, cost\\
RF drone detection & Long range & Expensive, emits or needs dedicated sensors; doesn't locate the casualty\\
Manual map and compass resection & Doctrine, no tech & Slow, needs training and calm\\
\bottomrule
\end{tabularx}
\end{table}

\paragraph{What is hard to copy.}
\begin{itemize}
  \item The ask-back loop: choosing the next question by expected error removed, with answer models for 14 question types.
  \item Radio-only operation end to end: callsign routing, local speech, replies approved by an operator.
  \item Data over the voice channel, with packets and a modem built for lossy radios and mirrored in the browser.
  \item A phone-sized sound model trained to hear through a radio. On simulated radio it catches drones that the big model misses.
  \item 3D line of sight from LiDAR, so ``I can't see the keep'' becomes evidence instead of noise.
\end{itemize}
Each piece alone is known technology. The combination, built around a speaker who doesn't know where he is, is ours.

\section{Go-to-market}

\nanofig{gtm}{Go-to-market: from a field trial at Vincennes to funded pilots, then fielding and a civilian version.}
{Horizontal roadmap timeline infographic, 16:9, plain white background, flat vector, navy line with cyan milestones. Four phases
left to right: ``Now: hackathon prototype, Vincennes'' (laptop icon), ``3 months: field trials with a military medical
training unit, real radios'' (radio icon), ``6–12 months: funded pilot (EDF consortium / AID / DIANA), Ukraine partner
trial'' (handshake icon), ``12–24 months: fielded site licences + civilian version for rescue services'' (map and headset
icon). Under the line, small grey labels ``validate'', ``prove'', ``scale''. Exactly these labels. Colour palette (use only these): navy \#0D1110, cyan \#1AA7C4, amber \#E09A2B, red \#E5484D, green \#3FA877, grey \#8B9A94, plus white for text on dark shapes.}

\begin{enumerate}
  \item \textbf{Prove it in the field (0--3 months).} A walk at Vincennes with GPS off; recordings through real tactical
        radios; a session with a military medical training unit to validate the medical rules.
  \item \textbf{Funded pilot (6--12 months).} A pilot with one medical unit or training centre, paid by a defence innovation
        programme; a trial with a Ukrainian partner, where jamming and drones are daily conditions.
  \item \textbf{Fielding (12--24 months).} Site licences and area packs; integration with the battle management system the
        customer already uses.
  \item \textbf{Civilian version.} The same engine for callers who can't give an address (lost hikers, foreign tourists,
        disaster zones without a network). In the EU, network-based caller location is missing in under 3 \% of 112 calls in
        most countries \cite{eu112}, but Advanced Mobile Location needs a SIM and coverage \cite{eu112}. Rescue in
        mountains and disaster areas is where both fail.
\end{enumerate}

\section{How we will validate the business}\label{sec:validate}

\begin{itemize}
  \item \textbf{Need:} interviews with combat medics and CCP coordinators on how often casualty positions are unknown or
        wrong, and how long locating takes today.
  \item \textbf{Fit:} timed exercises, the same scenario with and without Near: time to a usable position, error, number of
        transmissions.
  \item \textbf{Price:} what a medical service pays today for comparable software per site, and how area data is procured.
  \item \textbf{Channel:} whether customers buy directly, through an integrator, or through an innovation programme first.
\end{itemize}

\section{Risks}

\begin{table}[htbp]\centering\small
\caption{Main risks and mitigations.}\label{tab:risks}
\begin{tabularx}{\linewidth}{>{\raggedright}p{4cm}X}
\toprule
Risk & Mitigation\\
\midrule
Real-world accuracy lower than simulated & Field trials before any claim; honest radius; ask-back loop and compass bearings as fallback\\
Radio vocoders garble tones and sounds & Test on target radios; use radio data mode or mesh radio where available; voice-only mode always works\\
Medical-device regulation & Decision-support framing; early regulatory advice; medical validation of every rule\\
Export control and dual-use rules & Early legal check; friendly-forces-only design\\
Long defence sales cycles & Funded pilots, integrators, civilian version in parallel\\
Large incumbents add a similar feature & Speed, the combination of hard pieces, area-pack tooling, field relationships\\
Emissions detected & Few, short transmissions; offline phone; no extra emitter\\
\bottomrule
\end{tabularx}
\end{table}

\section{Team and ask}

\paragraph{Team Fricassé.} Taher Ayadi (pipeline, audio, radio and offline modes, models) and Aziz Maatoug
(geometry, localisation, 3D line of sight, casualty model). We built Near in one weekend at EDTH Paris.

\paragraph{Ask.}
\begin{itemize}
  \item Access to real tactical radios and a military medical training unit for a field trial.
  \item A medical mentor to validate thresholds, windows and team rules.
  \item Introductions to integrators and to defence innovation programmes for a funded pilot.
\end{itemize}
